\documentclass[11pt,a4paper]{article}

\usepackage[margin=2.3cm]{geometry}
\usepackage{amsmath,amssymb}
\usepackage{booktabs}
\usepackage{xcolor}
\usepackage[most]{tcolorbox}
\usepackage[colorlinks=true,linkcolor=blue!55!black,citecolor=blue!55!black,
            urlcolor=blue!55!black]{hyperref}

\newcommand{\dalpha}{\Delta\alpha}
\newcommand{\ReG}{\operatorname{Re}G_{\perp}}
\newcommand{\ee}{\mathrm{e}}
\newcommand{\ii}{\mathrm{i}}
\newcommand{\dd}{\mathrm{d}}
\newcommand{\unit}[1]{\,\mathrm{#1}}

\title{\bfseries Challenge 20: Torsional levitated optomechanics\\[2pt]
\large Derivation of the librational frequency $\omega_t$ and of the
light--induced torsion--torsion coupling $g$}
\author{}
\date{}

\title{Challenge 20: Problem}
\author{}
\begin{document}
\maketitle
\section{Statement and strategy}

Two identical dielectric prolate ellipsoids (semi-major axis $a$, two equal
semi-minor axes $b$, relative permittivity $\epsilon_r$, mass density $\rho$) are
levitated in two Gaussian optical tweezers that propagate along $\hat z$, are
linearly polarised along $\hat x$, carry power $P_0$, have waist $w_0$ and wave
vector $k$, and whose foci are separated by $R$ along $\hat x$.  The long axes
align with the polarisation and perform small librations about it.  We must
produce $\omega_t$ and $g$ in

\begin{equation}
  H=\hbar\omega_t\left(a_1^{\dagger}a_1+a_2^{\dagger}a_2\right)
   +\hbar g\left(a_1^{\dagger}a_2+a_1a_2^{\dagger}\right).
  \label{eq:H}
\end{equation}
\end{document}
